% arara: lualatex
% arara: bib2gls
% arara: lualatex
% arara: lualatex
\documentclass[toc=listof]{scrreport}

\usepackage{fontspec}
\setmainfont{Linux Libertine O}
\usepackage
 [
%   debug=showwrgloss
 ]{texjavahelp}

\hypersetup{colorlinks,linkcolor=blue}
\graphicspath{{images/}}

\TJHRequireIcons{texjavahelplibicons}{16}{24}

\newcommand{\Jmakepdfx}{Jmakepdfx}

\title{\Jmakepdfx}
\author[Nicola L.C. Talbot]{Nicola L.C. Talbot\\\href{https://www.dickimaw-books.com/}{\nolinkurl{dickimaw-books.com}}}

\TeXJavaHelpLoadResources{jmakepdfx,jmakepdfx-dict}

\begin{document}
\maketitle

\frontmatter
\tableofcontents

\mainmatter
\chapter{Introduction}
\label{sec:intro}

\Jmakepdfx\ is a Java \gls{GUI} front-end application that uses GhostScript to convert
PDF files to PDF/X compliant. This application comes with NO
WARRANTY. In order to use this application you need to have both
Java (at least JRE~8) and GhostScript installed. You will also need
the \TeX\ Java Help Library
(\href{https://github.com/nlct/texjavahelp}{\nolinkurl{github.com/nlct/texjavahelp}}). 

This application uses GhostScript's
\code{/prepress} setting, which is analogous to Acrobat Distiller's
\qt{Prepress Optimized}. This will likely create a larger PDF than
your original.

\begin{warning}
Note that PostScript does not support transparency. If your
PDF file includes transparent elements, the resulting PDF/X file
will end up as a raster image.
\end{warning}

\chapter{Installation}
\label{sec:install}

If \Jmakepdfx\ is available in your \TeX\ distribution, the best
method of installing it is via your \TeX\ distribution's package
manager. The instructions below are provided in the event that 
\Jmakepdfx\ is not available in your \TeX\ distribution.
To test if \Jmakepdfx\ is correctly installed, run the following in
a command prompt or terminal:
\begin{terminal}
\app{jmakepdfx} \switch{version}
\end{terminal}
This should print the version information.

If \Jmakepdfx\ is not available in your \TeX\ distribution, 
there are two methods of installation. The first is to install on
the TEXMF tree, where TEXMF denotes the distribution root, for
example \filefmt{/usr/local/texlive/2026/texmf-dist}
or \homefilefmt{texmf} (Unix-like systems) 
or \WindowsOSPath{c:/texlive/2026/texmf-dist} (Windows).

\begin{itemize}
\item TEXMF\filefmt{/scripts/jmakepdfx/jmakepdfx-helpset.tjh}
\item TEXMF\filefmt{/scripts/jmakepdfx/jmakepdfx.jar}
\item TEXMF\filefmt{/scripts/jmakepdfx/jmakepdfx.tlu}
\end{itemize}

The \file{texjavahelplib.jar} library also needs to be installed:
\begin{itemize}
\item TEXMF\filefmt{/scripts/texjavahelp/texjavahelplib.jar}
\end{itemize}
(There will be other files in the \filefmt{texjavahelp}
subdirectory. These are only required if building the
documentation. They are not needed to run \app{jmakepdfx}.)

For Unix-like systems, add a symbolic link to the
\TeX Lua script \file{jmakepdfx.tlu} called
\filefmt{jmakepdfx} somewhere on your system PATH.
This script locates \file{texjavahelplib.jar} and adds it to Java's
class path when running \file{jmakepdfx.jar} (the application jar
file).

On Windows, copy \filefmt{runscript.exe} to \filefmt{jmakepdfx.exe}.
This should find the \file{jmakepdfx.tlu} script and run it.

\begin{information}
The \ext+{tjh} file is a \TeX\ Java Help file, which contains the
files used by the in-application help that's used by the
\menu{help.manual} menu item and dialog \btn{help} buttons.
The localisation files used for messages and the textual parts of
\gls{GUI} elements are embedded in the \file{jmakepdfx.jar} file.
Translations can be added via pull request at
\href{https://github.com/nlct/jmakepdfx}{\nolinkurl{github.com/nlct/jmakepdfx}}.
\end{information}

The second approach is to place the \file{texjavahelplib.jar}
library in the same directory as the \file{jmakepdfx.jar}
application.
\begin{itemize}
\item \filefmt{jmakepdfx-helpset.tjh}
\item \filefmt{jmakepdfx.jar}
\item \filefmt{texjavahelplib.jar}
\end{itemize}
In this case, the \TeX Lua script is not required and the
\file{jmakepdfx.jar} application can be run as normal for your
system. For example, double-click on the \file{jmakepdfx.jar}
in your filer, if supported, or via the terminal:
\begin{terminal}
java -jar /path/to/jmakepdfx.jar
\end{terminal}
(Replace \filefmt{/path/to/jmakepdfx.jar} with the path to
\file{jmakepdfx.jar}.)

\chapter{Converting a PDF File}
\label{sec:conversion}

\FloatFig[hbtp]
 {fig:mainwindow1}
 {%
   \includeimg[alt={Main window (no files selected)}]{mainwindow1}
 }
 {Main Window (Initial View)}

The default setting is \gls{GUI} mode. (See \sectionref{sec:cmdline}
for command line invocation.) On startup, the main window is as
shown in \figureref{fig:mainwindow1}.

\menudef{menu.file.input}
Select the PDF file that needs conversion. This can be done via the 
\menu{file.input} menu item or the corresponding button on the
toolbar or you can use the \btn{widget.inputfile.choosein} button to the
side of the \widget{inputfile} field or you can drag and drop the
PDF from a filer window into the \widget{inputfile} field.

\begin{information}
If you startup \Jmakepdfx\ from the command line in \gls{GUI} mode
with \switch{in} then the specified file will automatically be
set in the \widget{inputfile} field.
\end{information}

Once the input file has be selected, \Jmakepdfx\ will run
GhostScript to obtain the author, title and page count.
These will be added to the \widget{pdfinfo.title} fields.

\widgetdef{pdfinfo.title}
The \widget{pdfinfo.title} fields consist of two editable fields
(for the title and author) and two non-editable fields (for the page
count and file size).

\widgetdef{pdfinfo.pdftitle}
The \widget{pdfinfo.pdftitle} field may be used to specify the PDF
title metadata. If found in the input file, this field will
automatically be set. You can edit this field. For example, if you
want to change the title, if it's incorrect, or add a title if one
wasn't found in the input file.

\widgetdef{pdfinfo.pdfauthor}
The \widget{pdfinfo.pdfauthor} field may be used to specify the PDF
author metadata. If found in the input file, this field will
automatically be set. You can edit this field. For example, if you
want to change the author, if it's incorrect, or add the author's
name if one wasn't found in the input file.

The \inlineglsdef{pdfinfo.pagecount} field shows the number of pages found
in the input file. 
The \inlineglsdef{pdfinfo.filesize} field shows the file size of the input file.
These fields are for informational purposes and are not editable.

\widgetdef{profile.title}
The profile may be set to either \inlineglsdef{profile.greyscale}
for grayscale or \inlineglsdef{profile.cmyk} for CMYK.

\widgetdef{profile.use_icc}
If you want to use an \gls{ICC} profile, ensure that the
\widget{profile.use_icc} checkbox is selected. You will need to
specify an \gls{ICC} file. This can be done in the
\dialog{properties} dialog,
which can be opened with the \menu{settings.editsettings} menu item.

\menudef{menu.file.output}
Select the output PDF file. This can be done via the 
\menu{file.output} menu item or the corresponding button on the
toolbar or you can use the \btn{widget.outputfile.chooseout} button to the
side of the \widget{outputfile} field. All these methods will open a
file chooser where you can navigate to the desired directory and
supply a filename. A default filename will be provided that's derived from
the input filename. Change it as applicable.

\begin{information}
If you startup \Jmakepdfx\ from the command line in \gls{GUI} mode
with \switch{output} then the specified file will automatically be
set in the \widget{outputfile} field.
\end{information}

If you have previously created an output file that you want to
overwrite, you can also drag and drop the
PDF from a filer window into the \widget{outputfile} field.

\menudef{menu.tools.convert}
The \menu{tools.convert} menu item or the corresponding button on the
main window is only enabled once both an input file and output file
have been specified (see \figureref{fig:mainwindow2}). 

Once you have selected all the appropriate settings,
select \menu{tools.convert} or click on 
\btn{menu.tools.convert} to run GhostScript.

\begin{information}
The resulting PDF file is likely to be bigger than the original.
\end{information}

You may be prompted to supply the path to GhostScript if 
\Jmakepdfx\ can't detect it. \Jmakepdfx\ searches for:
\filefmt{gs}, \filefmt{gswin64c}, \filefmt{gsos2}.
If you have a 32-bit Windows and require \filefmt{gswin32c.exe}
then you will need to set the path in the \dialog{properties} dialog
(see \sectionref{sec:settings}).

\begin{warning}
If your file size is very large, you may get a
\qt{\errmsg[\meta{n}]{timedout}} message, in which case  
increase the
\widget{properties.timeout} value in the \dialog{properties} dialog.
\end{warning}

\FloatFig
 {fig:mainwindow2}
 {%
   \includeimg[alt={Main window (input and output files selected)}]{mainwindow2}
 }
 {Main Window (Input and Output Files Selected)}

\widgetdef{button.auto_open}
If the \widget{auto_open} checkbox is selected and the conversion is
successful then \Jmakepdfx\ will attempt to open the PDF file.
Depending on your system settings, you may need to supply the path
to your preferred PDF viewer. Some systems may not allow this
action.

\begin{information}
If an error occurs, try running \app{jmakepdfx} with
\switch{verbose} (which will print the process call to STDOUT) 
and \switch{norm-tmp} (to retain the file used by the process).
There may be a file permission problem.
\end{information}

\chapter{Application Settings}
\label{sec:settings}

TODO

\chapter{Command Line Invocation}
\label{sec:cmdline}

\appdef{jmakepdfx}

\Jmakepdfx\ defaults to \gls{GUI} mode, but it can be run in batch
mode from the command line. The following switches are primarily
intended for batch mode use, but can also be used with the \gls{GUI}
mode, where they will populate or select the corresponding \gls{GUI} fields.

\switchdef{help}
Prints help information to the terminal and exits.

\switchdef{version}
Prints version information to the terminal and exits.
This information includes the version details for
the \TeX\ Java Help Library (\file{texjavahelplib.jar}).

\switchdef{gui}
Switches on the \gls{GUI} mode (default setting).

\switchdef{batch}
Switches on the batch mode setting (that is, switch off \gls{GUI}
mode).

\begin{important}
Note that even if \switch{batch} is present in the argument list, 
\file{jmakepdfx.tlu} will run with its own \gls{GUI} flag on.
\end{important}

\switchdef{nogui}
A synonym for \switch{batch}.

\switchdef{in}
Specifies the input PDF file. This switch is required for batch mode. With
\gls{GUI} mode, it will set the file in the input file field when the main
\Jmakepdfx\ window is opened.

The \switch{in} switch is optional. You can simply specify the
filename without \switch{in} (provided the filename doesn't start
with a hyphen).

\switchdef{output}
Specifies the output PDF file. This switch is required for batch mode. With
\gls{GUI} mode, it will set the file in the output file field when the main
\Jmakepdfx\ window is opened.

The \switch{output} switch is optional. You can simply specify the
filename without \switch{output} (provided the filename doesn't start
with a hyphen). If you omit \switch{output} you must specify the
output file \emph{after} you have specified the input file.

\switchdef{title}
Sets the PDF title for the output file metadata. If omitted, the title
obtained from the input file will be used, if set.

\switchdef{author}
Sets the PDF author for the output file metadata. If omitted, the author
obtained from the input file will be used, if set.

\switchdef{gray}
Use grayscale profile.

\switchdef{cmyk}
Use CMYK profile.

\switchdef{icc}
Indicates that the ICC file should be used: either the default ICC
file that has previously been saved via the \gls{GUI} application
settings, see \sectionref{sec:settings}, or the file specified with
\switch{icc-file} (if set).

\switchdef{noicc}
Indicates that the ICC file should not be used, even if one has been specified.

\switchdef{icc-file}
Specifies the path to the ICC file.

\switchdef{timeout}
Sets the maximum process time (in milliseconds).
If the GhostScript process runs on too long, the timeout will cancel
the process to prevent it from becoming a zombie process.
A large PDF file with a CMYK profile may take a while, in which case
you may need to extend the maximum process time.

\switchdef{rm-tmp}

On exit, removes the temporary PostScript file created by \app{jmakepdfx} that's 
provided for use with the GhostScript process. This setting is the
default.

\switchdef{norm-tmp}

Indicates that the temporary PostScript file created by \app{jmakepdfx}
should not be removed on exit. This setting is provided for
debugging if something goes wrong with the GhostScript process.

\begin{information}
If a problem occurs, try running \app{jmakepdfx} with
\switch{verbose} (which will print the process call to STDOUT) 
and \switch{norm-tmp} (to retain the file used by the process).
\end{information}

\switchdef{verbose}

Switches on verbose mode.

\switchdef{noverbose}

Switches off verbose mode.

\switchdef{debug}

Switches on debug mode.
The \meta{level} may be omitted. If present, it should be an integer
that indicates the debugging level. A value of 0 is equivalent to
\switch{nodebug}.

\switchdef{nodebug}

Switches off debug mode.

\printmain

\printindex

\end{document}
